\par\medskip
\noindent\textbf{Corollary \SGARefTarget{sga7.sga7i.evi.n6-3}{6.3}.}
Let $X_0$ be a $k$-scheme of finite type, and let $x\in X_0$ be a
non-degenerate quadratic singularity rational over $k$.

\smallskip
\noindent\textbf{(1)} Let $(R,\mathfrak X)$ and
$(R,\mathfrak Y)$ be two deformations of $X_0$, with
$R\in\operatorname{ob}(\widehat{\underline C}_\Lambda)$. Assume that
$k=k^2$. Then
\[
\mathfrak X_{(\widehat x)}\simeq
\mathfrak Y_{(\widehat x)}
\]
as $R$-adic formal schemes if and only if
\[
I_R(\mathfrak X_{(x)}|R)=I_R(\mathfrak Y_{(x)}|R).
\]

\smallskip
\noindent\textbf{(2)} Let
$R\in\operatorname{ob}(\widehat{\underline C}_\Lambda)$ be a regular
local ring, and let $(R,\mathfrak X)$ be a deformation of $X_0$.
Then $\mathcal O_{\mathfrak X,x}$ is regular if and only if
\[
I_R(\mathfrak X_{(x)}|R)\not\subset\mathfrak m_R^2;
\]
equivalently, $I_R(\mathfrak X_{(x)}|R)$ is a regular divisor in
$\operatorname{Spf}(R)$.

\smallskip
\noindent\textit{Proof.}
Let $(R_0,Z)$ be a formal versal deformation of
$\operatorname{Spec}(\widehat{\mathcal O}_{X_0,x})$. Since $x$ is a
non-degenerate quadratic singularity rational over $k$, Lemma \SGARefLink{sga7.sga7i.evi.n6-2}{6.2}
allows us to take
\[
R_0\simeq\Lambda[[t]],
\qquad
Z\simeq\operatorname{Spf}
 \bigl(R_0[[X_1,\ldots,X_n]]/(F-t)\bigr).
\]

For (1), the deformations
$(R,\mathfrak X_{(\widehat x)})$ and
$(R,\mathfrak Y_{(\widehat x)})$ are induced from $(R_0,Z)$ by
continuous local $\Lambda$-algebra homomorphisms
\[
\varphi,\psi:R_0\longrightarrow R.
\]
Put $a=\varphi(t)$ and $b=\psi(t)$. Then
\[
\begin{aligned}
\mathfrak X_{(\widehat x)}
&\simeq\operatorname{Spf}
 \bigl(R[[X_1,\ldots,X_n]]/(F-a)\bigr),\\
\mathfrak Y_{(\widehat x)}
&\simeq\operatorname{Spf}
 \bigl(R[[X_1,\ldots,X_n]]/(F-b)\bigr).
\end{aligned}
\]


Suppose that
$I_R(\mathfrak X_{(x)}|R)=I_R(\mathfrak Y_{(x)}|R)$; equivalently,
$(a)=(b)$ in $R$. Then $a=ub$ for some unit $u\in R$. We construct a matrix $T\in\operatorname{GL}_n(R)$ such that
$F(TX)=uF(X)$.\footnote{Editorial proof repair: the printed inference that
every unit of $R$ has a square root need not hold in characteristic $2$.
The following argument retains the original hypotheses.}
Write $\mathfrak m=\mathfrak m_R$, and interpret polynomial congruences
coefficientwise. Since $k=k^2$, choose $c\in R^\times$ whose residue
satisfies $\bar c^2=\bar u$, and put $T_1=cI_n$. Then
$F(T_1X)\equiv uF(X)\pmod{\mathfrak m}$.
By the nondegeneracy condition in Lemma \SGARefLink{sga7.sga7i.evi.n6-2}{6.2}, the coefficient matrix of
the linear forms $\partial F/\partial X_i$ is invertible modulo
$\mathfrak m$, hence over $R$.

Suppose $T_\nu\in\operatorname{GL}_n(R)$ satisfies
$F(T_\nu X)\equiv uF(X)\pmod{\mathfrak m^\nu}$, with $\nu\geq1$.
The forms $(\partial F/\partial X_i)(T_\nu X)$ form an $R$-basis of
the linear forms. Consequently, there are linear forms $L_i$ with
coefficients in $\mathfrak m^\nu$ such that
\[
E_\nu(X):=uF(X)-F(T_\nu X)
 =\sum_i L_i(X)\frac{\partial F}{\partial X_i}(T_\nu X).
\]
Indeed, every quadratic monomial is a product of two linear forms,
and expressing one factor in this basis preserves its coefficient ideal.
Define $H_\nu$ by $H_\nu X=(L_i(X))_i$, and set
$T_{\nu+1}=T_\nu+H_\nu$. The quadratic identity gives
\[
F(T_{\nu+1}X)=F(T_\nu X)+E_\nu(X)+F(H_\nu X)
 \equiv uF(X)\pmod{\mathfrak m^{\nu+1}},
\]
since the coefficients of $F(H_\nu X)$ lie in
$\mathfrak m^{2\nu}\subseteq\mathfrak m^{\nu+1}$.
Moreover, $T_{\nu+1}\equiv T_\nu\pmod{\mathfrak m^\nu}$.
Completeness and separatedness therefore give a limit
$T\equiv cI_n\pmod{\mathfrak m}$ with $F(TX)=uF(X)$; in particular,
$T$ is invertible. The substitution $X\mapsto TX$ sends
$F-a$ to $u(F-b)$ and induces the required isomorphism
\[
R[[X_1,\ldots,X_n]]/(F-a)
 \xrightarrow{\ \sim\ }
R[[X_1,\ldots,X_n]]/(F-b).
\]
The converse follows from invariance of $I_R$ under $R$-isomorphisms.

For (2), there is a continuous local $\Lambda$-algebra homomorphism
\[
\varphi:R_0=\Lambda[[t]]\longrightarrow R
\]
such that
\[
\mathfrak X_{(\widehat x)}
\simeq\operatorname{Spf}
 \bigl(R[[X_1,\ldots,X_n]]/(F-\varphi(t))\bigr).
\]
Moreover,
\[
I_R(\mathfrak X_{(x)}|R)
=I_R(\mathfrak X_{(\widehat x)}|R)
=(\varphi(t)),
\]
and
\[
R[[X_1,\ldots,X_n]]/(F-\varphi(t))
=\widehat{\mathcal O}_{\mathfrak X,x}
\]
is regular if and only if $\mathcal O_{\mathfrak X,x}$ is regular.
Since $R$ is regular,
$R[[X_1,\ldots,X_n]]$ is a regular local ring. If $M$ denotes its
maximal ideal, then its quotient by $F-\varphi(t)$ is regular if and
only if
\[
F-\varphi(t)\notin M^2,
\]
which is equivalent to
$\varphi(t)\notin\mathfrak m_R^2$, and hence to
\[
I_R(\mathfrak X_{(x)}|R)\not\subset\mathfrak m_R^2.
\]
\hfill$\square$
